\chapter{Agent（智能体）}

\section{概述}

Agent（智能体）是 OpenClaw 中一个完全独立作用域的"大脑"，拥有自己的：

\begin{itemize}
  \item \textbf{工作区}（文件、AGENTS.md/SOUL.md/USER.md、本地笔记、人设规则）
  \item \textbf{状态目录}（\texttt{agentDir}）：认证配置文件、模型注册表
  \item \textbf{会话存储}（聊天历史 + 路由状态）：\texttt{\textasciitilde{}/.openclaw/agents/<id>/sessions}
\end{itemize}

配置路径：\texttt{agents.list[]}

\section{单智能体模式（默认）}

不做任何配置时，OpenClaw 运行单个智能体：

\begin{itemize}
  \item \texttt{agentId} 默认为 \texttt{main}
  \item 工作区默认为 \texttt{\textasciitilde{}/.openclaw/workspace}
  \item 状态默认为 \texttt{\textasciitilde{}/.openclaw/agents/main/agent}
\end{itemize}

\section{多智能体模式}

Gateway 网关可托管多个智能体并行运行，每个智能体拥有：

\begin{itemize}
  \item 不同的工作区和人格（AGENTS.md、SOUL.md）
  \item 独立的认证和会话
  \item 独立的 Skills
\end{itemize}

\subsection{添加智能体}

\begin{lstlisting}[language=bash]
openclaw agents add work
openclaw agents list --bindings
\end{lstlisting}

\subsection{配置示例}

\begin{lstlisting}[language=json]
{
  agents: {
    list: [
      {
        id: "home",
        default: true,
        name: "Home",
        workspace: "~/.openclaw/workspace-home"
      },
      {
        id: "work",
        name: "Work",
        workspace: "~/.openclaw/workspace-work"
      }
    ]
  }
}
\end{lstlisting}

\section{工作区文件}

工作区是智能体的家目录，包含以下标准文件：

\begin{center}
\begin{tabular}{ll}
\toprule
\textbf{文件} & \textbf{用途} \\
\midrule
\texttt{AGENTS.md} & 操作指南（每个会话开始时加载） \\
\texttt{SOUL.md} & 人设、语气和边界 \\
\texttt{USER.md} & 用户偏好 \\
\texttt{memory/} & 记忆文件目录 \\
\texttt{skills/} & 每智能体 Skills \\
\bottomrule
\end{tabular}
\end{center}

\section{智能体超时}

\begin{itemize}
  \item 默认超时：\texttt{agents.defaults.timeoutSeconds}（默认 600 秒）
  \item \texttt{agent.wait} 超时：默认 30 秒（仅等待，不停止智能体）
\end{itemize}

\section{常用命令}

\begin{lstlisting}[language=bash]
# Run agent directly (no inbound chat)
openclaw agent --to +15555550123 --message "status update"
openclaw agent --agent ops --message "Summarize logs"

# List agents
openclaw agents list --bindings
\end{lstlisting}
